\section{Principal charts and marked component orbits on the lower model}

% BEGIN C76P1
\paragraph{C76P1. A principal period atlas on each pinned lower row.}
Let $Q:\mathrm{Par}(D)$ satisfy $Q.\mathrm{Pins}$ and let
$\mathrm{InR}(D,j)$. Put $D_w=Q.Dw(j)$, $X_o=Q.Xo(j)$, $c=Q.c$,
$M=Q.l.M$, $E=D.\mathrm{host}.E$, $O=\mathcal O_E$ and $r=D.R.r$.
Use the inclusion $\sigma_0:c.F\hookrightarrow\mathbb C$.
There is a single $U':\mathrm{Unf}(X_o,c,\sigma_0,r)$ satisfying
$\mathrm{IsUnf}$. Its nonempty finite component set has a common
Hermitian matrix $G$ and scalar $\xi=-\xi_H$, where
$\xi_H=D.\mathrm{host}.\xi_0$. Every component has
$\mathrm{HasPrinc}(X_k,M)$ and $\chi_k(\lambda,\gamma)=1$ for
every weight and group argument. Here $X_k=U'.X(k)$.
Each $X_k$ extends to an original $\mathrm{DataX}$ with
\[
 L'_k=\mathfrak A^{-1}\Lambda_k,\qquad
 \Lambda_k=\operatorname{span}_O\{a u:a\in\mathfrak A,\ u\in L'_k\},
 \qquad \mathfrak A=D.\mathrm{host}.\mathrm{AA}\ne0 .
\]

\begin{proof}
The row signature pins give $p+q=r\ge1$. The level pins give
$l.A=1$, hence $l.M'=\operatorname{lcm}(M,1)=M$.
The inclusion is the good embedding for the given coefficient
record. To check the sign in the period theorem, write $s^2=-d$
for the host generator and $N$ for the positive norm in PCL8.
Then $\xi_H=-1/(Ns)$ and $\bar\xi_H=-\xi_H$.
At a host-type embedding $s$ maps to $i\sqrt d$, by the root and
embedding pins. Thus $\operatorname{Im}\tau(-\xi_H)<0$.
Apply PCL1 to $D_w$ with multiplier $-\xi=\xi_H$.
It gives an $O$-linear injection $\phi:D_w.\mathrm{La}\to E^r$
whose image spans, and a nondegenerate $G=\bar G^{\mathsf T}$
with $D_w.\mathrm{EF}(a,b)=
\operatorname{Tr}_{E/\mathbb Q}(-\xi h_G(\phi a,\phi b))$.
No choice of an integral basis is needed.

D74U2 supplies both section laws of NFL22 on this $X_o$.
Apply NFL22 with the just chosen $G,\xi,\phi$. Its flag models,
normalized framed period points and test presentations therefore
refer to this lower coefficient field and this trace pairing.
Its group condition reads $(\gamma-I)v=Mw$ with $w\in\Lambda_k$
for every $v\in\Lambda_k$, for a unitary lattice stabilizer.
Adding $v$ gives exactly the congruence in the component definition
of $\mathrm{HasPrinc}$. The characters and the parameters on every
component are the specified ones of NFL22. In addition, NFL22 gives
the test presentation for each framed point and the Fourier exponent
match at every cusp rank since $l.A=1$.

For the lattice assertion, the integral period identification
$\phi_k$ in $\mathrm{IsUnf}$ carries a unimodular alternating
lattice onto $\Lambda_k$. Its trace form is therefore integral
and perfect, by PCL6. By PCL8,
$(\xi)\mathfrak A\bar{\mathfrak A}\mathfrak D_E=O$ and
$|\tau\xi|=\xi R(D)\in(0,1]$ for every embedding. Apply PCL5
to this component with $g=M\ge3$. The Hermitian dual and the
normalization fields are those of PCL5, and inversion of the
fractional ideal gives the asserted product on the literal lattice.
This construction precedes the choice of a section or weight.
\end{proof}
% END C76P1

% BEGIN C76P2
\paragraph{C76P2. Transport of the weighted bound to these charts.}
Let $U$ be the given distinguished row uniformization of the same
$X_o,c,\sigma_0$, satisfying $\mathrm{IsUnf}$. For a dominant
$(a,n)$ and $f\in\mathrm{sec}(X_o,c,a,n)$ suppose the readings
on $U$ have component bounds $C_k$. For a component $k'$ of $U'$
and a CM embedding $b$ put
\[
 h_k(b)=2\pi|\mathrm{relab}(c,\sigma_0,b)(\xi_k)|,\qquad
 h'_{k'}(b)=2\pi|\mathrm{relab}(c,\sigma_0,b)(\xi'_{k'})|.
\]
A bound on the new reading is
\[
 \sum_{k\in U.\mathrm{Comp}}|C_k|
 \prod_b\prod_{0\le i<D_w.\mathrm{mult}(b)}
       \left(\frac{h'_{k'}(b)}{h_k(b)}\right)^{a(b,i)/2}.
\]
In particular this applies to every power of a cusp form occurring in
ICL10 at its resulting dominant weight.

\begin{proof}
Use upper triangular Gram--Schmidt inside the prescribed signed flag.
The transformation law for the leading Gram minors cancels the Borel
weight in the pointwise weighted norm. At the resulting orthonormal
flag the absolute values of these minors are one. It therefore
suffices to bound the scalar reading on orthonormal flags.

On such a flag the invariant Hermitian form of the period torus
has diagonal entries $2|\mathrm{relab}(b)\xi_k|$ in the ordered
tangent coordinates. Indeed, in the positive block the period
coordinates are Hermitian column pairings, while in the
complementary block they are the corresponding inverse-matrix
coordinates. Using the orthogonal columns in these two formulas,
the imaginary part of the diagonal form is the alternating trace
pairing on every pair of periods. Periods span the real tangent
space. A Hermitian form is determined by its imaginary part there
by evaluation also on $iv$; positivity on each coordinate fixes
the sign of the purely imaginary embedding of $\xi_k$.
The Gaussian metric multiplies this matrix by $\pi$. The normalized
theta generator has squared mean norm one by the fundamental-domain
normalization in the framed chart. The same calculation applies on
$U'$.

Take a new orthonormal frame and represent its framed moduli point
by an old orthonormal flag followed by a Borel change. This uses
the framed covering of $U$ and right triangular orthonormalization.
The new and the old frames with that change give the same point,
so have the same metric there. Indeed a rigidified metric on the
theta line with the prescribed curvature is unique: the logarithm
of the ratio is harmonic on the compact connected torus, hence
constant, and the indicated normalization fixes the constant. The
induced tangent form is intrinsic to the polarization. Consequently
the section-change scalar has absolute value one. On the $b$ block
the triangular change matrix has diagonal entries of absolute value
$\sqrt{h'_{k'}(b)/h_k(b)}$: the scalar Gram identity makes the
matrix, after this common rescaling, unitary, and a triangular
unitary matrix is diagonal. The algebraic weight rule on framed
points gives exactly the displayed product. Only the old component
containing the chosen frame contributes to this pointwise estimate;
summing the nonnegative bounds over the finite component set gives
a uniform bound. The new reading already has its holomorphy and
automorphy from $\mathrm{IsUnf}$. For cusp forms the original bounds
follow from the cusp condition, and multiplication retains the cusp
condition in the weight used in ICL10.
\end{proof}
% END C76P2

% BEGIN C76P3
\paragraph{C76P3. The alternating form and the dual half-marking.}
Fix a reference component of C76P1. Its domain and alternating form
can be used for all components. Write $A=D_w.\mathrm{La}$,
$I=\mathrm{Mfr}(l)$, $\bar I=l.M_b$, so that $I\bar I=(M)$
and $I+\bar I=O$. Let $L$ be a component lattice and
$m:A\to L$ its actual period marking. With the primitive root
$\sigma_0(\zeta_M(c))=\exp(2\pi i t/M)$, there is an
$O$-linear map $\delta:A\to L$ satisfying
\[
  B(\delta x,m a)\equiv -t\,D_w.\mathrm{EF}(x,a)\pmod M
         \qquad(x\in A,\ a\in IA).                         \tag{C76.1}
\]
It induces an isomorphism $A/IA\to L/IL$. Conversely its reduction
determines the original marking on $IA/MA$.

\begin{proof}
The CM type and positive-sign index of each chart are the ones in
$\mathrm{IsUnf}$ for $D_w$. The period and model conventions give
$\kappa_0=2\xi$ in their common coordinates. To see the sign, the
periods of a frame span the complex tangent space over $\mathbb R$.
The imaginary alternating form of its positive period Hermitian
matrix is determined on those periods. If $\kappa_0=-2\xi$,
the positive form obtained from the signed flag coordinates and the
positive polarization form would have opposite imaginary parts on
every pair of periods. Their sum would have zero imaginary part
on the real tangent space. Evaluating also on $(v,iv)$ would make
that sum zero, impossible in the positive dimension $r|\Psi|$.
Thus the native alternating form $B=\mathrm{EF}(X)$ agrees with
the unimodular component trace form of C76P1 on the whole lattice.
The same argument identifies the reference and component domains,
up to the literal reindexing of embeddings.

The integral period isomorphism transports the given half-marking
isomorphism to $m:IA/MA\simeq IL/ML$. The alternating perfect
pairings on $A/MA$ and $L/ML$ pair the two complementary halves
perfectly with one another. Indeed $O/(M)=O/I\times O/\bar I$,
and moving an idempotent across either pairing conjugates it, thereby
interchanging the factors. Each factor is isotropic and the perfect
pairing identifies it with the additive dual of the opposite factor.
Equation (C76.1) therefore defines a unique map on the complementary
quotients $A/IA\to L/IL$: it is $-t$ times the adjoint of the
inverse of the given marking. Here $t$ is invertible modulo $M$.
Moving $O$-scalars across the pairings twice proves $O$-linearity.
Since $A$ is finite projective over the Dedekind ring $O$, lift its
composite with $A\to A/IA$ through $L\to L/IL$. This gives the
required integral $\delta$. In the exponential form of the original
definition the argument is $c_{\rm sep}a$; it has the same reduction
as $a$ since $c_{\rm sep}-1\in\bar I$. This proves the sign and
root convention in $\mathrm{IsDual}$. Conversely the same perfect
pairing and invertibility of $t$ recover $m|_{IA/MA}$ from
$\delta\bmod I$. Unitary transport preserves this relation,
including its root.
\end{proof}
% END C76P3

% BEGIN C76P4
\paragraph{C76P4. A single tuple and pairwise distinct marked lattice orbits.}
Let $n\ge r$, and use an original classical level $\mathcal D$ tied
to the identity cover of the lower row. Write $\mathfrak m=IJ$
for its prescribed coprime factorization, $\beta_i$ for the
$r$ level vectors, $\Delta=\mathcal D.\mathrm{grp}$ and $X$ for the
reference component. There exist a single admissible $n$-tuple
$x^0$ in $A$, rational unitary matrices $a_k$, and transported
dual markings $a_k\delta_k$ such that
\[
 L_k=a_k\Lambda_k\in\mathcal D.\mathrm{lats},\qquad
 a_k\delta_k(x^0_i)\equiv\beta_i\pmod I\quad(0\le i<r).
\]
The last $n-r$ entries of $x^0$ vanish. The congruence is understood
in the reference lattice at primes dividing $I$. If
$g\in\Delta$ and $gL_k=L_{k'}$, then $k=k'$.

\begin{proof}
Use C76P3 on each component. In the reference component lift
the residues of $\beta_i$ through $\delta_0$ to $b_i\in A$.
Since the $\beta_i$ span modulo $\mathfrak m$ they generate
modulo $I$. The $b_i$ thus generate $A/IA$. Set
$x^0=(b_0,\ldots,b_{r-1},0,\ldots,0)$. This gives
$\mathrm{AdmX}$ for every allowed lower-row threshold: its defining
span includes the first $r$ positions. For each component the
vectors $\delta_k b_i$ generate $\Lambda_k/I\Lambda_k$.
Choose $r$ generators modulo $J$ of $\Lambda_k$: locally it is
free of rank $r$, and the Artinian quotient modulo $J$ is a
finite product of local rings. For $J=O$ the condition is empty.
Combine the two lists by the Chinese remainder theorem to get
$u_i\in\Lambda_k$ generating modulo $\mathfrak m$ and reducing
to $\delta_k b_i$ modulo $I$. In the general source formula the
first list has a factor $d.\mathrm{pd}$; this factor equals one
on the identity cover used here.

Here is why all the congruences can be realized by a rational
\emph{unitary} matrix. At primes $P\mid\mathfrak m$ the $u_i$ and
the reference $\beta_i$ are bases of their respective lattices by
Nakayama. Write $B_0 U_0^{-1}$ for the matrix between the two lists
of columns. The level condition
$\mathfrak m+\bar{\mathfrak m}=O$ selects at most one of each
conjugate pair of primes. Above the corresponding primes of
$E^+$ the local unitary group is $\mathrm{GL}_r(E_P)$ on this
chosen factor; the conjugate factor is determined by the Hermitian
adjoint. Rational unitary elements approximate these local elements
to any fixed precision. For clarity, the Cayley transform
$S\mapsto (1-S)(1+S)^{-1}$ identifies the open skew-adjoint
matrices with $1+S$ invertible with the corresponding open of
the unitary group. In the split local general linear groups,
perturb each prescribed matrix inside its congruence neighborhood
to have no eigenvalue $-1$ (multiply by a scalar arbitrarily near
one avoiding the finite exceptional set). The inverse transforms
are in the local skew-adjoint spaces. The rational skew-adjoint
space is a vector space over $E^+$; weak approximation in a basis
meets all their neighborhoods. The nonzero local determinants
ensure its rational Cayley transform is defined. Taking the
precision high enough preserves the desired congruence on the
finitely many specified columns, even before clearing their
denominators.

Choose this transform as $a_k$. Then $L_k=a_k\Lambda_k$ equals
the reference lattice locally at $P\mid\mathfrak m$, since the
residue lists are bases. Both lattices are self-dual for the
alternating trace pairing. Local duality gives equality also
at the conjugate primes. Clearing the finitely many remaining
local denominators in the two inclusions gives one common multiplier
prime to $\mathfrak m\bar{\mathfrak m}$, exactly the comparison
in $\mathcal D.\mathrm{lats}$. The carried marking and dual
marking on $L_k$ are those of C76P3, with the stated residues.

Suppose $gL_k=L_{k'}$. The level condition on $g$ fixes every
$\beta_i$ modulo $I$. On the $I$-local lattices, which have just
been identified, we obtain $g a_k\delta_k b_i\equiv
a_{k'}\delta_{k'} b_i\pmod{I}$. Since the $b_i$ generate modulo
$I$, this holds on every element of $A$. It is an equality in
$L_{k'}/I L_{k'}$, since membership in $I L_{k'}$ imposes only
these extra local conditions on a vector already in $L_{k'}$.
Recovering the marking by (C76.1) proves
$g a_k m_k(a)\equiv a_{k'}m_{k'}(a)\pmod{M L_{k'}}$ for all
$a\in IA$. The other marking has level one. Thus
$a_{k'}^{-1}ga_k$ gives a marked integral isometry of the two
original period data. PUN23 identifies the resulting marked period
points after transporting a period by this isometry; HUF14 gives
disjoint images for distinct marked indices. In the actual atlas
of C76P1 the period index is the component index
($\mathrm{unfOf}$ uses these indices directly). Hence $k=k'$.
\end{proof}
% END C76P4

% BEGIN C76P5
\paragraph{C76P5. The classical level attached to the identity cover.}
For the lower identity cover of ICL11, a level $\mathcal D$ satisfying
every field of $\mathrm{LevTie}$ in C76P4 exists on the reference
chart of C76P1. Its scalar in the tying identity equals one.

\begin{proof}
Take $\mathfrak m=I=\mathrm{Mfr}(l)$ and $J=J_d=O$.
The level fields give $I\bar I=(M)$ and $I+\bar I=O$.
The product is nonzero and proper since $M\ge3$. As conjugation
preserves $O$, both factors are nonzero proper ideals. The integer
$M$ belongs to $I$ and is odd, so $I+(2)=O$. The $O$-lattice of
the reference chart is projective of rank $r$; its reduction modulo
$I$ is free of rank $r$ on each of the finitely many local factors
of $O/I$. Combine bases of these factors and lift their $r$ vectors
to $\beta_0,\ldots,\beta_{r-1}$ in the original lattice.
Take the residue subgroup $\{1\}\subset\mathrm{GL}_r(O/I)$
and its trivial character $\lambda$. Take the prescribed integral
$w=H.w_0=(1+s)/2$, so $w+\bar w=1$.

The two classical moduli required in $\mathrm{LevTie}$ now coincide
with $I=IO$, and $I$ is coprime to $O$. A residue matrix differs
entrywise from $1$ by residues of elements of $I$ precisely when
it is $1$. This proves the prescribed subgroup and determinant
conditions. To check the last field, take any admitted lattice,
map $\delta$, tuple, and matrix $y$ satisfying the conditions
there. On the first $r$ rows $y_i\equiv\beta_i\pmod I$
in the localized reference lattice. Hence the reduction predicate
$\mathcal D.\mathrm{Red}$ holds with the integral identity matrix.
On the remaining rows the required congruence is $y_i\equiv0\pmod I$.
The native finite function $\mathrm{datF}(y)$ therefore equals
$\lambda(1)=1$. The other side is the identity-cover function
$\mathrm{ph}=1$ on the zero auxiliary quotient of ICL11.
This gives the tying identity for every such $y$ with a single
scalar $1$ even before choosing $\delta$ and the lattice. Thus
C76P4 applies on the actual lower level.
\end{proof}
% END C76P5

% BEGIN C76P6
\paragraph{C76P6. The letter and weight dictionary on the reference chart.}
Use the letter enumeration $e$ from the period witness in C76P1.
Every dominant $(a,n)$ has an original classical dictionary
$\mathrm{Dic}$ on the reference component with this same enumeration.

\begin{proof}
Write $q=X.p_0=D_w.q$, $p=D_w.p$, and $g=|\Psi|r$.
Self-duality in C76P3 supplies a symplectic $\mathbb Z$-basis of
the reference lattice: pair a primitive vector with a vector of
pairing one, split their unimodular plane, and repeat on its
orthogonal complement. Choose the sign of the second half as in
the original alternating matrix. Since the rank is $2g$, this
gives the required $g+g$ basis.

The good relabeling commutes with conjugation and identifies the
host type with $\Psi$. If $\tau$ belongs to the host type and
$\tilde\tau=\mathrm{relab}(c,\sigma_0,\tau)$, send the letters
$(\tau,i)$, $0\le i<p$, to $(\tilde\tau,q+i)$, and
$(\bar\tau,i)$, $0\le i<q$, to $(\tilde\tau,q-1-i)$.
These disjoint lists exhaust the $r$ indices at every
$\tilde\tau$. Composing the permutation with $e$ gives
$\mathrm{idx}$; adjoining the unchanged copy index $b<n$
gives the polynomial letter permutation. Substitution into the
period definition gives respectively the inverse-frame coordinate
and the Hermitian column coordinate. In particular, differentiation
in a renamed letter is differentiation in the original tangent unit
vector in copy $b$ and coordinate $e^{-1}(\tau,i)$
(or $e^{-1}(\bar\tau,i)$ for the second list).

For each of the two weakly decreasing lists $a_i$, take the
successive differences $a_i-a_{i+1}$ with final entry $a_{d-1}$
for a list of length $d>0$. The tail sums recover $a_i$.
The resulting Borel exponents in positions $q+i$ and $q-1-i$
are $a(\tau,i)+\lfloor n/2\rfloor$ and
$-a(\bar\tau,i)-\lfloor n/2\rfloor$, as in $\mathrm{wtX}$.
Empty blocks contribute no letters or exponents. Under the
polynomial permutation, the defining row homogeneities, highest
row operations and paired harmonic contractions are the same
equations, term by term. In the fixed multidegrees their coefficients
are in $c.F$; kernels of finite matrices commute with field extension.
This gives the membership and spanning fields of $\mathrm{Dic}$.
It also preserves a positive final exponent in either indicated
nonempty list whenever required, by the definition of its final
difference. ICL15 supplies the further integral spanning assertion
for the chosen analytic lattice.
\end{proof}
% END C76P6

% BEGIN C76P7
\paragraph{C76P7. Every component stabilizer fixes the normalized vacuum.}
Let $n$ be even, $L=L_k$ a moved lattice of C76P4--C76P5,
and $g\in\Delta$ satisfy $gL=L$. On the native vacuum and
multiplier of the classical level, one has
\[
             (g-I)L\subset ML,\qquad \mu(g,L)=1 .
\]
For a bounded chart reading $P_k$, the transported function
$x\mapsto C_k^{-1}P_k(a_k^{-1}x)$ on the reference domain
is a bounded stabilizer-invariant form of the weight in C76P6
for every nonzero comparison scalar $C_k$.

\begin{proof}
The elements of $\Delta$ and their inverses act as identity on the
reference lattice modulo $I$ at its support. Indeed these conditions
are a group, and are the defining residue conditions on generators
because $\beta_i$ give bases there. Since $L$ is the reference
lattice at these places, $g,g^{-1}$ act as identity there modulo
$I$. On the conjugate component of each local alternating perfect
pair, the action is the inverse adjoint of that on the first.
Hence $g$ is identity also modulo $\bar I$ there. At all other
places $M$ is a unit. Since $g$ stabilizes $L$, these assertions
give $(g-I)L\subset ML$ globally. The endomorphism
$(g-I)/M$ preserves $\mathcal D.\mathrm{sh}(L)$ as well:
the latter is $\{x:\bar{\mathfrak m}x\subset\mathfrak m L\}$,
and an $O$-linear endomorphism preserving $L$ preserves
$\mathfrak m L$ and commutes with these products. The shadow lattice is self-dual and is
stabilized by $g$, by the classical level laws.

It remains to fix the scalar on the vacuum line of the shadow.
A self-dual alternating CM lattice has an integral symplectic
basis adapted to the host action modulo eight. In that basis
the matrix of $g^{-1}$ on periods has off-diagonal blocks divisible
by eight and its indicated diagonal block has determinant $1$
modulo eight by NFL4. Its hypotheses hold here: transport the
integral alternating form on the shadow to a CM datum, using its
rational isometry with the reference space, unimodularity, and
the reference signature. The inverse is congruent to identity
modulo $M\ge3$ on this very lattice. In particular the determinant
of this CM-linear isometry is one: it is an integral unit of
absolute value one at every embedding, and its difference from one
is divisible by $M$. If nonzero, the absolute norm of this
difference would be at once at least $M^{[E:\mathbb Q]}$ and
at most $2^{[E:\mathbb Q]}$.

The signed vacuum in period coordinates is a product of $n$ copies
of the same theta section, with its Gaussian factors and the
normalizing determinant to the power $n/2$. The multiplier in this expression is independent of the period
frame. To see this, transport the finite test by the Weil
operator of the fixed symplectic change of basis. The diagonal
Heisenberg generators force its support into one coset of the
integer matrices, with all rows shifted by the same first
half-characteristic. Integral unit translations in each coordinate
specify every value on that coset from one value by the second
half-characteristic. Both characteristics are obtained from the
semicharacter on the fixed basis vectors. Thus the transported
test is one nonzero constant times this signed coset indicator,
before choosing a frame. Its theta sum factors by rows in every
frame. Gaussian completion of squares gives the product just
described with this same constant. The change of period normalization
contributes the same half-characteristic and multiplier square
$\det j$ by NFL2--NFL5 (the determinant congruence makes the
residual quadratic character one). If the one-variable theta
factors after the Gaussian change are compared up to a scalar
$z$, the product law and the existence of a point where the theta
section is nonzero permit division in each variable, and
$z^2=\det j$. The $n$ copies thus contribute
$(\det j)^{n/2}$, which cancels the change of the normalizing
determinant. For $n=0$ the empty product has scalar one directly.
The action fixes the vacuum as a vector. Its defining transformation
law reads $g v_L=\mu(g,L)\lambda(1)v_L$, whence $\mu(g,L)=1$
because $\lambda(1)=1$ and $v_L\ne0$.

Conjugation by $a_k$ now gives
$a_k^{-1}g a_k\in\Gamma_k$ by C76P1. Its character is one.
The pullback is therefore invariant under the whole stabilizer.
Left unitary translation is holomorphic, preserves all the leading
Gram minors and commutes with the right Borel action. The literal
reindexing of embeddings in C76P6 identifies the Borel characters
on the two charts. Hence the pullback has the same weighted bound
up to $|C_k|^{-1}$. Because C76P4 gives disjoint component orbits
and the vacuum multiplier is one on stabilizers, the extensions
along these orbits by the prescribed multiplier law agree for
all choices of a mover. This is exactly the stabilizer and orbit
input of ICL16.
\end{proof}
% END C76P7
